\section{Introduction}

Educational institutions increasingly deliver course content through browser-based learning platforms. These have matured for text, video and 2D-interactive material, but integrating interactive 3D visualization remains an unresolved engineering problem, despite documented pedagogical value~\cite{Li2021,Yang2024} and established technology pipelines~\cite{Rodrguez2021,Hensen2023}.

The deployment context imposes constraints that are not typical of general-purpose 3D rendering systems. Students access content using heterogeneous hardware, including older laptops with integrated graphics, mid-range gaming laptops, and laboratory workstations of varying age, over which the institution has no control. In addition, the installation of native software is typically prohibited on managed devices, so the content must run in a standard browser tab without any plugins. Content development teams consist of subject-matter experts and engineers without graphics-programming experience, meaning that the development of each scenario requires engineering effort in addition to the work involved in creating the educational content itself.

Existing solutions expose a trade-off between development convenience and deployment efficiency~\cite{Evans2014}. For educational institutions the practical requirement is therefore a lightweight browser runtime that preserves a high-level development model, integrates into an existing platform, and can be self-hosted.

This paper presents the design and implementation of the 3D runtime of the KSU24 virtual learning environment at Kherson State University. Implemented in TypeScript and distributed as an npm library, the engine provides a Unity-inspired component model, including \texttt{GameObject}, \texttt{Component}, \texttt{Transform}, \texttt{Material}, and lifecycle hooks such as \texttt{awake}, \texttt{start}, and \texttt{update}, while fully encapsulating the underlying Three.js backend. This allows content developers to create browser-based educational 3D scenarios using familiar high-level abstractions without interacting directly with the low-level graphics library. Ten educational scenarios have been built on it and packaged for the platform. The work further investigates three optimization strategies targeting distinct runtime costs: GPU-compressed texture delivery using KTX2/Basis Universal for reducing texture memory consumption and first-render cost, shader pre-compilation for reducing rendering initialization latency, and dirty-flag transform propagation for reducing per-frame transform-processing overhead. The evaluation covered 87 configurations across three benchmark scenes, using three graphics adapters on two machines, with each optimization assessed against the metric it is intended to affect.